\documentclass[10pt,a4paper]{amsart}
\usepackage[margin=19mm]{geometry}
\usepackage{amsmath,amssymb,mathtools}
\usepackage[scr=rsfs,cal=euler]{mathalfa}
\usepackage{xfrac}
\usepackage[unicode,hidelinks,bookmarks=false]{hyperref}
\newcommand{\ie}{i.e.\ }
\newcommand{\cf}{cf.\ }
\newcommand{\resp}{respectively\ }
\newcommand{\Subsection}{\subsection}
\providecommand{\nobreakdash}{\nobreak\hbox{-}}
\setlength{\emergencystretch}{2em}
\multlinegap0pt
\usepackage{amssymb}
\usepackage{dsfont}
\usepackage{algorithm}\usepackage{algpseudocode}
\usepackage{mathtools}
\usepackage{tikz}
\usepackage[shortlabels]{enumitem}
\usepackage{bm}
\usepackage{tcolorbox}
\newtheorem{lemma}{Lemma}
\title{Text of the title}
\author{Michael Herty, Hui Huang, Hicham Kouhkouh}
\date{Source: arXiv:2608.23066v1 (CC BY 4.0)}
\begin{document}
\maketitle

\noindent\emph{Source and licence.} Text of the title. Version arXiv:2608.23066v1 (CC BY 4.0), distributed by arXiv under the Creative Commons Attribution 4.0 International licence, http://creativecommons.org/licenses/by/4.0/, as recorded in the arXiv licence metadata for this exact version and shown on the abstract page https://arxiv.org/abs/2608.23066v1. Full licence notice: \texttt{licenses/arxiv-2608.23066-CC-BY-4.0.txt}.\par\smallskip

\noindent\textit{Editorial scope.} Counted unit: the article's lemma lemma (source0.tex lines 1386--1388). The article's complete original proof of it is delivered with it (source0.tex lines 1390--1495, one \texttt{proof} environment of 106 lines). No mathematics is written, repaired or re-derived: the statement and the proof are the article's own lines, in the article's own notation, and no retained line is substituted or re-flowed.  Retained uncounted as the source-local context the proof needs: lines 609--611; lines 613--615.  Nothing was omitted from any retained span, so the package carries no omission record.  No retained line contains a citation, so the wrapper carries no bibliography and no bibliography entry.  No retained line contains a figure environment, an image-inclusion command or drawing code, so no image is delivered and the package declares no asset dependency.  Editorial formatting only, in addition: an AMS \texttt{amsart} preamble that restates the article's own declarations and macros used by the retained lines, namely the environments \texttt{lemma} on separate counters, together with the standard packages those lines need.  The retained environments print in this document as Lemma 1 in order of appearance, whereas the article prints the same items with the numbering of its own full document.  The complete native source of the article as distributed in the arXiv e-print for this exact version is bundled unchanged as \texttt{sources/208388/source0.tex}.
        \begin{equation}\label{eq:osl_condition}
            \langle \nabla \mathscr{E}(x) - \nabla \mathscr{E}(y), x - y \rangle_H \ge -C_{\mathscr{E}} \|x - y\|_H^2.
        \end{equation}

        \begin{equation}\label{eq:quadratic_condition}
            \mathscr{E}(y) - \mathscr{E}(x) \le \langle \nabla \mathscr{E}(x), y - x \rangle_H + \frac{L_{\mathscr{E}}}{2} \|y - x\|_H^2.
        \end{equation}

\begin{lemma}
Let $\mathscr{E}:H\to\mathbb R$ be Fr\'echet differentiable. If $\mathscr{E}$ is simultaneously $C$-semiconvex and $L$-semiconcave, then $\mathscr{E}\in C^{1,1}(H)$,  and $\operatorname{Lip}(\nabla\mathscr{E}) \le L+C$. In particular, the conclusion remains true if $\mathscr{E}$ satisfies \eqref{eq:osl_condition}-\eqref{eq:quadratic_condition}.
\end{lemma}

\begin{proof}
Recall a function is said to be $L$-smooth, if its gradient is $L$-Lipschitz. 

\textbf{Step 1.} \textit{(Convexity and semiconcavity imply $L$-smoothness.)}\\
We first show that if a Fr\'echet differentiable function
$\Phi:H\to\mathbb R$ is convex and $M$-semiconcave, then $\|\nabla\Phi(x)-\nabla\Phi(y)\|_H \le M\|x-y\|_H$ for all $x,y\in H.$

If $M=0$, then $\Phi$ is both convex and concave, hence affine, and there is nothing to prove. Assume therefore that $M>0$. Since subtracting a linear functional preserves both convexity and $M$-semiconcavity, the function
\[
\Psi(z):=\Phi(z)-\langle\nabla\Phi(y),z\rangle_H
\]
is convex and $M$-semiconcave. Moreover, $\nabla\Psi(y)=0$,  so  $y$ is a global minimizer of $\Psi$. By $M$-semiconcavity,
\[
\Psi(v) \le \Psi(u) + \langle\nabla\Psi(u),v-u\rangle_H + \frac{M}{2}\|v-u\|_H^2, \qquad u,v\in H.
\]
Indeed, this is the first-order characterization of the concavity of
$\Psi-\frac{M}{2}\|\cdot\|_H^2$. Applying this inequality with
\[
u=x, \qquad v=x-\frac1M\nabla\Psi(x),
\]
yields
\[
\begin{aligned}
    \Psi(v)
    & \le \Psi(x) - \frac{1}{M}\|\nabla\Psi(x)\|_H^2 + \frac{M}{2} \left\| \frac1M\nabla\Psi(x) \right\|_H^2 \\
    & = \Psi(x) - \frac{1}{2M} \|\nabla\Psi(x)\|_H^2.
\end{aligned}
\]
Since $y$ minimizes $\Psi$, $\Psi(y)\le\Psi(v)$, and therefore
\[
\Psi(y) \le \Psi(x) - \frac1{2M} \|\nabla\Psi(x)\|_H^2.
\]
Recalling the definition of $\Psi$, we have $\nabla \Psi(x) = \nabla \Phi(x) - \nabla \Phi(y)$, and using the previous inequality we obtain
\[
\Phi(x)-\Phi(y) - \langle\nabla\Phi(y),x-y\rangle_H \ge \frac{1}{2M} \|\nabla\Phi(x)-\nabla\Phi(y)\|_H^2.
\]
Interchanging the roles of $x$ and $y$ and adding the two inequalities
gives
\[
\langle \nabla\Phi(x)-\nabla\Phi(y), x-y \rangle_H \ge \frac{1}{M} \|\nabla\Phi(x)-\nabla\Phi(y)\|_H^2.
\]
Finally, the Cauchy--Schwarz inequality implies
\[
\|\nabla\Phi(x)-\nabla\Phi(y)\|_H^2 \le M \|\nabla\Phi(x)-\nabla\Phi(y)\|_H \|x-y\|_H,
\]
and hence
\[
\|\nabla\Phi(x)-\nabla\Phi(y)\|_H \le M\|x-y\|_H.
\]

\textbf{Step 2.} \textit{(Semiconvex and semiconcave imply $C^{1,1}$.)}\\
Define
\[
\mathscr{F}(x) := \mathscr{E}(x)+\frac{C}{2}\|x\|_H^2, \qquad \mathscr{G}(x) := \frac{L}{2}\|x\|_H^2-\mathscr{E}(x).
\]
By the $C$-semiconvexity and $L$-semiconcavity of $\mathscr{E}$, both $\mathscr{F}$ and $\mathscr{G}$ are convex. Moreover,
\[
\mathscr{F}+\mathscr{G} = \frac{L+C}{2}\|\cdot\|_H^2.
\]
Consequently, $\mathscr{F}$ is $(L+C)$-semiconcave, since
\[
\mathscr{F}(x)-\frac{L+C}{2}\|x\|_H^2 = -\mathscr{G}(x)
\]
is concave. Similarly,
\[
\mathscr{G}(x)-\frac{L+C}{2}\|x\|_H^2 = -\mathscr{F}(x),
\]
so $\mathscr{G}$ is also $(L+C)$-semiconcave.

If $L+C=0$, then necessarily $L=C=0$. In this case, $\mathscr{E}$ is both convex and concave, hence affine, and $\nabla\mathscr{E}$ is constant. We may therefore assume that $L+C>0$.

Applying Step 1 to $\mathscr{F}$ and $\mathscr{G}$, with $M=L+C$, yields for all $x,y\in H$
\[
\|\nabla\mathscr{F}(x)-\nabla\mathscr{F}(y)\|_H \le (L+C)\|x-y\|_H
\quad \text{ and } \quad 
\|\nabla\mathscr{G}(x)-\nabla\mathscr{G}(y)\|_H \le (L+C)\|x-y\|_H.
\]

Let
\[
p:=\nabla\mathscr{E}(x)-\nabla\mathscr{E}(y), \qquad h:=x-y.
\]
Then
\[
\nabla\mathscr{F}(x)-\nabla\mathscr{F}(y) = p+Ch, \qquad \text{ and } \qquad  \nabla\mathscr{G}(x)-\nabla\mathscr{G}(y) = Lh-p.
\]
Since
\[
p = \frac{L}{L+C}(p+Ch) - \frac{C}{L+C}(Lh-p),
\]
the triangle inequality gives
\[
\begin{aligned}
    \|p\|_H
    & \le \frac{L}{L+C}\|p+Ch\|_H + \frac{C}{L+C}\|Lh-p\|_H \\
    & = \frac{L}{L+C}\|\nabla\mathscr{F}(x)-\nabla\mathscr{F}(y)\|_H + \frac{C}{L+C}\|\nabla\mathscr{G}(x)-\nabla\mathscr{G}(y)\|_H \\
    & \le \frac{L}{L+C}(L+C)\|h\|_H + \frac{C}{L+C}(L+C)\|h\|_H 
     = (L+C)\|h\|_H.
\end{aligned}
\]
Hence,
\[
\|\nabla\mathscr{E}(x)-\nabla\mathscr{E}(y)\|_H \le \bigl(L+C_{\mathscr{E}}\bigr)\|x-y\|_H, \qquad x,y\in H.
\]
Therefore, $\nabla\mathscr{E}$ is globally Lipschitz continuous and $\mathscr{E}\in C^{1,1}(H)$. 
\end{proof}

\end{document}
